\documentclass[10pt,a4paper]{amsart}
\usepackage[margin=19mm]{geometry}
\usepackage{amsmath,amssymb,mathtools}
\usepackage[scr=rsfs,cal=euler]{mathalfa}
\usepackage{xfrac}
\usepackage[unicode,hidelinks,bookmarks=false]{hyperref}
\newcommand{\ie}{i.e.\ }
\newcommand{\cf}{cf.\ }
\newcommand{\resp}{respectively\ }
\newcommand{\Subsection}{\subsection}
\providecommand{\nobreakdash}{\nobreak\hbox{-}}
\setlength{\emergencystretch}{2em}
\multlinegap0pt
\usepackage{amssymb}
\usepackage{xcolor}
\usepackage{mathtools}
\usepackage[shortlabels]{enumitem}
\newtheorem{lem}{Lemma}
\newcommand{\R}{{\mathbb R}}
\newcommand{\lemref}[1]{Lemma~\ref{#1}}
\title{Ramaswami Type translation formulae for the polylogarithm functions}
\author{Pawan Singh Mehta, Biswajyoti Saha}
\date{Source: arXiv:2601.02921v1 (CC BY 4.0)}
\begin{document}
\maketitle

\noindent\emph{Source and licence.} Ramaswami Type translation formulae for the polylogarithm functions. Version arXiv:2601.02921v1 (CC BY 4.0), distributed by arXiv under the Creative Commons Attribution 4.0 International licence, http://creativecommons.org/licenses/by/4.0/, as recorded in the arXiv licence metadata for this exact version and shown on the abstract page https://arxiv.org/abs/2601.02921v1. Full licence notice: \texttt{licenses/arxiv-2601.02921-CC-BY-4.0.txt}.\par\smallskip

\noindent\textit{Editorial scope.} Counted unit: the article's lem sum (source0.tex lines 422--427). The article's complete original proof of it is delivered with it (source0.tex lines 428--456, one \texttt{proof} environment of 29 lines). No mathematics is written, repaired or re-derived: the statement and the proof are the article's own lines, in the article's own notation, and no retained line is substituted or re-flowed.  No further source-local context is needed: every cross-reference of every retained line resolves inside the two delivered spans.  Nothing was omitted from any retained span, so the package carries no omission record.  No retained line contains a citation, so the wrapper carries no bibliography and no bibliography entry.  No retained line contains a figure environment, an image-inclusion command or drawing code, so no image is delivered and the package declares no asset dependency.  Editorial formatting only, in addition: an AMS \texttt{amsart} preamble that restates the article's own declarations and macros used by the retained lines, namely the environments \texttt{lem} on separate counters, together with the standard packages those lines need.  The retained environments print in this document as Lemma 1 in order of appearance, whereas the article prints the same items with the numbering of its own full document.  The complete native source of the article as distributed in the arXiv e-print for this exact version is bundled unchanged as \texttt{sources/192194/source0.tex}.
 \begin{lem}\label{sum}
Let $q\geq2$ be an integer and $z$ be a primitive $q$-th root of unity. Then for every positive integers $k,m$, we have
\begin{align}\label{m-power-sum-z}
\sum_{h=1}^{k-1}z^{h}h^m =\frac{1}{q}\sum_{j=1}^{q-1} \frac{z(z^{j}-1)}{z-1}\left(E_{q,m}(j)-z^{k-1}E_{q,m}(k+j-1)\right).
\end{align}
\end{lem}

\begin{proof}
 Note that the polynomial $E_{q,n}(x)$ satisfies the formula,
\begin{equation}\label{diff-sum}
E_{q,n}(x)+E_{q,n}(x+1)+\cdots+E_{q,n}(x+q-1)=qx^n, \text{\ \ \ for all $n\geq 0$ and $x\in \R$.}
\end{equation}
This follows from the fact that
$$
\sum_{n\ge 0}\left(E_{q,n}(x)+E_{q,n}(x+1)+\cdots+E_{q,n}(x+q-1)\right)\frac{t^n}{n!}=qe^{xt}.
$$
By comparing the coefficients of $t^n$ on both the sides of the above equation, we get the identity \eqref{diff-sum}. 
%Note that for $n\geq 1$ if we substitute $x=0$ in \eqref{diff-sum} then we get that
%$$
%E_{q,n}(0)+E_{q,n}(1)+\cdots+E_{q,n}(q-1)=0,
%$$
Therefore we have,
\begin{align*}
q\sum_{h=1}^{k-1}z^{h}h^m&=\sum_{h=1}^{k-1}z^h\sum_{i=0}^{q-1}E_{q,m}(h+i)=\sum_{j=1}^{k+q-2}E_{q,m}(j)
\sum_{\substack{h+i=j, \\1\leq h\leq k-1,\\ 0\leq i\leq q-1}}z^h\\                                         
&=\sum_{j=1}^{q-1}E_{q,m}(j)\sum_{h=1}^{j}z^h+\sum_{j=q}^{k-1}E_{q,m}(j)\sum_{h=j-q+1}^{j}z^h+ \sum_{j=k}^{k+q-2}E_{q,m}(j)\sum_{h=j-q+1}^{k-1}z^h.
\end{align*}
Since $1+z+\cdots+z^{q-1}=0$, the second sum on the right-hand side of the above equation is zero. Hence we can write,
\begin{align*}                                            
 q\sum_{h=1}^{k-1}z^{h}h^m&=\sum_{j=1}^{q-1}E_{q,m}(j)z\frac{z^{j}-1}{z-1}+\sum_{j=0}^{q-2}E_{q,m}(k+j)z^{k-q}\sum_{h=j+1}^{q-1}z^h\\
 &=\sum_{j=1}^{q-1}E_{q,m}(j)z\frac{z^{j}-1}{z-1}+\sum_{j=0}^{q-2}E_{q,m}(k+j)z^{k+j+1}\frac{z^{q-j-1}-1}{z-1}\\
 %&=\sum_{j=1}^{q-1}E_{q,m}(j)z\frac{z^{j}-1}{z-1}+\sum_{j=0}^{q-2}z^k\frac{1-z^{j+1}}{z-1}E_{q,m}(k+j)\\
 &=\sum_{j=1}^{q-1}E_{q,m}(j)z\frac{z^{j}-1}{z-1}+\sum_{j=1}^{q-1}z^k\frac{1-z^{j}}{z-1}E_{q,m}(k+j-1).
\end{align*}
This completes the proof of \lemref{sum}.
\end{proof}

\end{document}
